\section{Evaluation of entire original Laurent-pulled module sections}

\begin{definition}\label{def:fc-scheme-module-laurent-evaluation}
\lean{EnvelopeProperty.SchemeModuleLaurentEvaluation.evaluate,EnvelopeProperty.SchemeModuleLaurentEvaluation.evaluate_coeff,EnvelopeProperty.SchemeModuleLaurentEvaluation.evaluate_monomial}
\leanok
\uses{thm:fc-scheme-module-laurent-coefficients,thm:fc-scheme-module-laurent-coefficient-formulas,thm:laurent-module-actual-evaluation}
Let $R$ be a commutative ring, $p:X\to\operatorname{Spec}R$ compact
and separated, and $M$ a quasicoherent module. Write $g:X_S\to X$
and $q:X_S\to\operatorname{Spec}S$ for the actual base-change projections,
where $S=R[T,T^{-1}]$. For any commutative ring $B$ that is an
$R$-module and any $R$-linear map $\phi:\Gamma(X,M)\to B$, there is
an evaluation function on the entire original pulled-module section
space into $B[T,T^{-1}]$. If $(m_i)$ is the unique original coefficient
family of $\sigma$, its evaluated polynomial has coefficient
$\phi(m_i)$ at every integer $i$. In particular it sends
$q^*(T^i)\eta_g(m)$ to $\phi(m)T^i$.
\end{definition}
\begin{proof}
Send the original section through the inverse genuine whole H0
base-change equivalence. Apply the preceding entire Laurent tensor
evaluation to this actual tensor. The geometric coefficient map is
exactly the original tensor coefficient map after that inverse, so
the coefficient identity follows at every integer. The original
monomial-unit section is the whole H0 image of $T^i\otimes m$;
the two inverse identities and the tensor monomial formula give
the stated evaluation. This constructs the evaluation from the
original linear map; no section-expansion or compatibility premise
is required.
\end{proof}

\begin{theorem}\label{thm:fc-scheme-module-laurent-gauss-maximum}
\lean{EnvelopeProperty.SchemeModuleLaurentEvaluation.laurentGauss_evaluate}
\leanok
\uses{def:fc-scheme-module-laurent-evaluation,thm:laurent-module-actual-gauss-maximum}
For any valuation $w$ of $B$ and any positive radius $r$, the actual
Laurent Gauss value of the evaluated original section is
\[
 \max_{i\in\operatorname{supp}(m_\bullet)}w(\phi(m_i))r^i.
\]
The support is that of the unique original module coefficients,
including coefficients killed by $\phi$ and all negative integer
weights. The empty maximum is zero.
\end{theorem}
\begin{proof}
Apply the entire tensor Laurent Gauss formula to the original section's
inverse H0 tensor. By construction its original tensor coefficients
are precisely the geometric coefficients of the original section,
and its evaluated polynomial is the section evaluation just defined.
Substituting these two identities into that formula proves the stated
maximum, including the zero section and coefficients killed by the
original linear map.
\end{proof}

\begin{theorem}\label{thm:fc-scheme-module-laurent-gauss-finite-family}
\lean{EnvelopeProperty.SchemeModuleLaurentEvaluation.laurentGauss_finite_family}
\leanok
\uses{thm:fc-scheme-module-laurent-gauss-maximum,def:fc-scheme-module-laurent-finite-coefficient-family,thm:laurent-module-actual-finite-family}
Every finite original family $G$ of sections of $g^*M$ determines the
finite coefficient and integer-weight family constructed above.
This family is fixed before $\phi$, $w$ and $r$. For every such choice,
the maximum over $G$ of the actual evaluated Laurent Gauss values is
exactly the weighted maximum $\max w(\phi(m))r^i$ over that same
original pair family $(m,i)$.
\end{theorem}
\begin{proof}
Take the image of the original finite section family under the inverse
whole H0 equivalence and apply the original tensor finite-family Gauss
formula. A finite maximum over this image is exactly the maximum over
the original section family, since the evaluation is the composition
with that inverse. The tensor coefficient-pair family equals the
constructed geometric pair family by the exact original
coefficient-family identity. Substitution gives the asserted equality.
The geometric pair family depends only on $G$, so the same family
works for all linear maps, valuations and positive radii.
\end{proof}
